\documentclass[12pt]{article}
\usepackage[utf8]{inputenc}
\usepackage[T1]{fontenc}
\usepackage{hyperref}
\title{Yali, Angguruk Language Resources for K-12 Curriculum Development}
\author{Compiled by Claude Sonnet 5.5 (claude-sonnet-5-5) for Andr\'as Kornai}
\date{}
\begin{document}
\maketitle

\section{Introduction}
Angguruk Yali (ISO 639-3: yli) is a Trans-New Guinea language of the Angguruk area, Yahukimo Regency, Papua, Indonesia. Most speakers also speak some Indonesian. The language is written in a Latin alphabet developed in the 1960s by German missionaries. The sources I found disagree on the date of the dictionary (2017 or 2021).

\section{Primary Educational Resources}
No primers, textbooks or school materials in Yali turned up in the search. The reference description says the language is not known to be taught in schools.

\section{Linguistic and Reference Materials}
Riesberg, Z\"ollner and Z\"ollner's \emph{A Yali (Angguruk) -- English -- German Dictionary} (Asia-Pacific Linguistics, ANU) has more than 3,000 headwords. Most have definitions and example sentences.
\begin{itemize}
\item HAL record: \url{https://halshs.archives-ouvertes.fr/halshs-03436264}
\item Dictionary introduction (PDF): \url{https://halshs.archives-ouvertes.fr/halshs-03466987v1/file/Yali_Dict_Intro_2021_FINAL.pdf}
\item ANU repository, ``Cultural advice'': \url{https://openresearch-repository.anu.edu.au/items/5162468a-0842-4f8d-91ca-2531122e95e3}
\item Glottolog reference record: \url{https://glottolog.org/resource/reference/id/666988}
\item Wikipedia, Yali language: \url{https://en.wikipedia.org/wiki/Yali_language}
\end{itemize}

\section{Religious and Community Resources}
Omniglot says Siegfried Z\"ollner, one of the missionaries, translated the Bible into Yali. The search did not confirm the publisher or current availability of the Bible.
\begin{itemize}
\item Omniglot, Yali writing: \url{https://www.omniglot.com/writing/yali.htm}
\item Joshua Project, Angguruk Yali in Indonesia: \url{https://joshuaproject.net/people_groups/15964}
\end{itemize}

\section{Digital and Audio Resources}
No audio, video or software resources turned up in the search. The dictionary is available online through HAL and ANU.

\section{Curriculum Development Context}
Of the four languages in this batch, Angguruk Yali has the best lexical foundation: an open dictionary with examples and an established Latin orthography. Reading materials, textbooks and audio would still have to be created. Before any programme relies on the Bible or the dictionary dates, those details should be checked.

\end{document}